% ============================================================================
% MODULE 6: LONG-TERM VARIABILITY
% ============================================================================
% Duration: 60 minutes (climate risk, yield variability, decision rules)
% Learning Goals: Assess yield risk over multi-year timescales; make robust decisions
% Prerequisites: Module 3, Module 5

\chapter{Module 6: Long-Term Variability}
\label{ch:module6}

\section{Learning Objectives}

By the end of this module, you will:
\begin{itemize}
  \item Analyze yield distribution over 20--50 year climate records
  \item Calculate risk metrics (coefficient of variation, probability of failure)
  \item Understand the role of rainfall variability in yield
  \item Design robust management strategies (that work in most years)
  \item Use multi-year simulations to assess profit variability
\end{itemize}

\section{Rainfall Variability}

\subsection{Inter-Annual Variability in Queensland}

Dalby (subtropical Queensland) experiences high year-to-year rainfall variability:

\begin{itemize}
  \item \textbf{Annual rainfall range}: Roughly 400--900 mm/year (long-term mean $\approx$ 600 mm)
  \item \textbf{Coefficient of Variation (CV) of rainfall}: 25--35\% --- meaning a ``typical'' year can be 250 mm above or below average
  \item \textbf{ENSO influence}: El Ni\~{n}o years tend to be drier; La Ni\~{n}a years wetter
  \item \textbf{In-season variability}: Drought can occur in any month, including during grain fill
\end{itemize}

\subsection{Why This Matters for Farming Decisions}

If a farmer optimises their management for the average year, they will:
\begin{itemize}
  \item Over-invest in N in drought years (money wasted; yield limited by water, not N)
  \item Under-invest in N in wet years (yield potential not reached)
\end{itemize}

The solution is \emph{risk-aware} management: choose strategies that are profitable across a range of seasons, not just the average.

% ---- IMAGE PLACEHOLDER ----
\begin{figure}[H]
  \centering
  \fbox{\parbox{0.85\textwidth}{\centering\vspace{1.5cm}
    \textit{[IMAGE: Bar chart of annual rainfall at Dalby, 1950--1969 from the weather file.}\\
    \textit{Annotate: long-term mean as horizontal dashed line; highlight 3 driest and 3 wettest years.]}
  \vspace{1.5cm}}}
  \caption{Annual rainfall variability at Dalby, 1950--1969 (from Exercise D weather data). Note the high year-to-year variation.}
  \label{fig:rainfall-variability}
\end{figure}
% ---- END PLACEHOLDER ----

\section{Yield Variability}

\subsection{Exercise D: 20-Year Variability Analysis}

Run simulation using 20 consecutive years from climate record (e.g., 1950--1969).

Expected outcomes:
\begin{itemize}
  \item Wet years: High rainfall during sowing window and grain fill $\to$ high ESW $\to$ high yield
  \item Dry years: Low or poorly-timed rainfall $\to$ water stress during grain fill $\to$ low yield
  \item Yield CV: 20--40\% is typical for rainfed wheat in subtropical Queensland
\end{itemize}

% ---- IMAGE PLACEHOLDER ----
\begin{figure}[H]
  \centering
  \fbox{\parbox{0.85\textwidth}{\centering\vspace{1.5cm}
    \textit{[IMAGE: Time series bar chart of annual wheat yield 1950--1969 from Exercise D.}\\
    \textit{Annotate: mean yield as dashed horizontal line; label the 3 lowest-yield years.]}
  \vspace{1.5cm}}}
  \caption{Annual wheat yield variability from the 20-year simulation (Exercise D). The dashed line shows the mean yield.}
  \label{fig:yield-variability}
\end{figure}
% ---- END PLACEHOLDER ----

\section{Risk Metrics}

\subsection{Coefficient of Variation}

\[
\text{CV} = \frac{\text{Std Dev}(\text{Yield})}{\text{Mean}(\text{Yield})}
\]

Higher CV = greater uncertainty; farmers prefer lower CV.

\subsection{Probability of Failure}

\[
P(\text{Yield} < Y_{\text{threshold}}) = \frac{\text{Number of years with Yield} < Y_{\text{threshold}}}{n \text{ total years}}
\]

Example: If 5 out of 20 years had yield $<$ 1500 kg/ha, then P(failure) = 5/20 = 25\%.

\subsection{Exceedance Probability (Cumulative Distribution)}

An exceedance curve shows the probability of exceeding a given yield value. It is constructed by:

\begin{enumerate}
  \item Sorting all 20 annual yields from lowest to highest
  \item Calculating rank: each yield gets a rank $r$ from 1 (lowest) to 20 (highest)
  \item Exceedance probability = $(n - r + 1) / n \times 100$
\end{enumerate}

\begin{keyconceptbox}
\textbf{Reading the exceedance curve:}
\begin{itemize}
  \item \textbf{P50} (50\% exceedance): The yield exceeded in 50\% of years --- the median. A better measure of typical performance than the mean.
  \item \textbf{P80} (80\% exceedance): The yield exceeded in 80\% of years --- the ``reliable'' yield.
  \item \textbf{P20} (20\% exceedance): The yield exceeded only 20\% of years --- the high-yield potential.
\end{itemize}
When advising farmers, P80 is the most practically useful: ``In 4 out of 5 years, you can expect at least X kg/ha.''
\end{keyconceptbox}

\section{Robust Decision Rules}

A \emph{robust} management strategy is one that performs reasonably well across the full range of likely seasons --- not just the average or best case.

\subsection{Example: N Rate Decision}

\begin{center}
\begin{tabular}{l|r|r|r|r}
\textbf{N rate (kg/ha)} & \textbf{Mean profit} & \textbf{Worst-year profit} & \textbf{CV (profit)} & \textbf{P(loss)} \\
\hline
0   & \$540 & \$120 & 45\% & 5\% \\
80  & \$730 & \$200 & 40\% & 10\% \\
160 & \$880 & \$180 & 48\% & 15\% \\
240 & \$850 & \$30  & 60\% & 25\% \\
\end{tabular}
\end{center}

(Indicative values; your Exercise D results will vary.)

\begin{itemize}
  \item 160 kg/ha maximises mean profit but has higher variability
  \item 80 kg/ha offers lower mean profit but more reliable income (lower CV)
  \item 240 kg/ha has the worst risk profile (high N cost, low payoff in dry years)
\end{itemize}

\begin{keyconceptbox}
\textbf{The risk aversion principle:} For subsistence or small-scale farmers, minimising catastrophic loss (P(loss)) is often more important than maximising mean profit. Use APSIM multi-year simulations to reveal these trade-offs explicitly before recommending a management strategy.
\end{keyconceptbox}

\section{Climate Change and Yield Risk}

A brief introduction (expanded in Module 7):

\begin{itemize}
  \item If mean temperature rises by 1--2\textdegree C, GDD accumulation speeds up --- crop matures faster, potentially shortening grain fill
  \item If rainfall declines (likely in some regions under SSP projections), yield CV will increase and P(failure) rises
  \item APSIM allows direct testing: modify the weather file (increase T$_{max}$, decrease rainfall) and re-run 20-year simulations to quantify impact
\end{itemize}

\begin{warningbox}
\textbf{Climate change increases variability, not just means.} Even if average rainfall stays the same, more extreme droughts and floods can increase yield CV and push P(failure) higher. This is why long-term simulations are essential --- single-season models miss this risk.
\end{warningbox}

\section{Summary}

\begin{keyconceptbox}
\textbf{Module 6 key points:}
\begin{enumerate}
  \item Dalby rainfall CV is 25--35\%; single-season analysis is insufficient for reliable management recommendations
  \item Run 20+ year simulations to capture the full range of climate outcomes
  \item Key risk metrics: CV (relative variability), P(failure) = probability yield falls below threshold, P50/P80 from exceedance curve
  \item Robust strategies have lower CV and P(failure), even if mean profit is slightly lower
  \item Climate change increases variability in addition to shifting mean conditions
\end{enumerate}
\end{keyconceptbox}

\section{What's Next?}

Module 7 concludes with climate scenarios and future outlook.

% ============================================================================
% END MODULE 6
% ============================================================================
